\color{unimain}{\subsection{Управление выключателем}}
\color{black}

\subsubsection{Общие сведения}

\begin{enumerate}[label=\arabic{section}.\arabic{subsection}.\arabic{subsubsection}.\arabic*,  align=left, labelsep=4pt, leftmargin=0pt, itemindent=57pt]

\item
Функциональный блок управления выключателем предназначен для подачи сигнала на включение
силового выключателя. 

\item
В случае присоединения линии к системе через два выключателя для каждого из них можно
задействовать свою логику управления выключателем. Далее приведено описание логики управления
выключателем В1, логика управления вторым выключателем аналогична.

\item
Команда на включение выключателя \textit{«Включить В1»} формируется:
\begin{itemize}
\item от команды оперативного включения выключателя (сигнал \textit{«РКВ В1(ЛВ)ВН»});
\item при срабатывании АПВ (сигнал \textit{«Ср. АПВ В1(ЛВ)ВН»}); 
\item от внешнего сигнала включения \textit{«Внеш. вкл. В1(ЛВ)ВН»}.
\end{itemize}

\begin{figure}[H]
\centering
\includegraphics[width=0.8\textwidth,height=0.8\textheight,keepaspectratio]{\fbpath/761.3391 Управление В/img/OP_Q.pdf}
\captionof{figure}{Функциональная схема логики отключения выключателя}\label{pic:OP_Q}
\end{figure}

\item
Оперативное включение выключателя осуществляется с контролем наличия разрешающего
сигнала от функционального блока КСН (принцип работы которого описан в пункте {\color{red}2.30} настоящего
руководства по эксплуатации). Максимальное время, в течение которого продолжается ожидание
условий оперативного включения, регулируется уставкой \textbf{«Tожид усл. опер. вкл. В1»}. По истечению
данной выдержки времени команда оперативного включения сбрасывается. Выдержка времени на
возврат \textbf{«Тпрод опер. вкл. В1»} задает время продления сигнала оперативного включения.

\item
При необходимости может быть введен подхват команды на включение от токового реле,
установленного в цепи ЭМВ, в этом случае сигнал на включение удерживается в сработанном состоянии
до исчезновения тока через ЭМВ (пока блок-контакты выключателя не разорвут цепь включения). Ввод
осуществляется программным ключом \textbf{«Подхват вкл. В1 от РТ ЭМВ»}. 

\item
Предусмотрена возможность блокировки многократных включений на КЗ. Если при наличии
сигнала включения будет протекать ток по электромагнитам отключения (что соответствует включению
на КЗ), то сигнал \textit{«Включить В1»} заблокируется, пока не пропадут сигналы включения (от АПВ или
оперативного включения). Блокировка вводится программным ключом \textbf{«Блок. от многокр. вкл. на КЗ
В1»}.

\item
При наличии сигнала \textit{«Блок. включения В1»}, формируемого функциональным блоком контроля
силового выключателя, включение выключателя блокируется.

\small
\begin{longtable}{|>{\centering\arraybackslash}m{5.3cm}|>{\centering\arraybackslash}m{4.2cm}|>{\centering\arraybackslash}m{1cm}|>{\centering\arraybackslash}m{3.3cm}|>{\centering\arraybackslash}m{1.8cm}|}
% Первая страница: полный заголовок
\caption{Перечень уставочных параметров УВ\hfill\vspace{-0.5\baselineskip}}\label{desc:tbl1}\\
\hline
\rowcolor{gray!30}
Наименование уставки & Диапазон & Шаг & Значение по умолчанию & Единица измерения \\
\hline
\endfirsthead

% Промежуточные страницы: «Продолжение таблицы»
\caption*{\hspace{3pt}\emph{Продолжение таблицы \ref{desc:tbl1}}\hfill\ }\vspace{-0.5\baselineskip} \\
\hline
\rowcolor{gray!30}
Наименование уставки & Диапазон & Шаг & Значение по умолчанию & Единица измерения \\
\hline
\endhead

%===m>ДТЗ|-
%===m
\end{longtable}
\normalsize	

\end{enumerate}